%% it/sections/05_npc.tex
%% ============================================================
%% 05_npc.tex — NPC storici per fazione
%% ============================================================

\chapter{I Personaggi Non Giocanti}
\markboth{I Personaggi Non Giocanti}{I Personaggi Non Giocanti}

\lettrine[lines=3]{I}{} personaggi non giocanti di \textit{Vespri 1282}
sono, dove possibile, figure storicamente attestate. Vengono presentati
con una scheda funzionale al gioco e una nota storiografica che aiuta
il GM a interpretarli con coerenza. La nota storiografica non è
un vincolo: è uno strumento.

\nota{Per tutti i GM}{
  Gli NPC storici sono punti di ancoraggio, non prigioni. Potete
  amplificarli, ridurli, spostarli. L'unica cosa che vi chiedete è:
  \textit{questa scelta è coerente con quello che sappiamo di questa
  persona?} Se la risposta è sì, andate avanti.
}

\secsep

%% ============================================================
%% NPC TRASVERSALI (disponibili a più tavoli)
%% ============================================================
\section{NPC Trasversali}

\subsection{Giovanni da Procida}
\textit{Medico, giurista, diplomatico. L'architetto della cospirazione — o la sua leggenda.}

\statblock{Giovanni da Procida}{Cospiratore Principale}%
{3}{4}{1}{3}

\textbf{Disponibile a:} Congiurati (presenza diretta), Baroni (emissario),
Clero (contatto indiretto)

\textbf{Ruolo in gioco:} Da Procida non è a Palermo — è il punto di
riferimento lontano della cospirazione. Può comparire attraverso lettere,
emissari, o in una scena speciale se il GM lo ritiene opportuno. La sua
presenza è più potente quando è assente: i personaggi discutono di lui,
interpretano le sue istruzioni, si chiedono se si fida di loro.

\textbf{Carattere:} Freddo, metodico, capace di affetto genuino solo
verso la causa. Ha visto troppo per fidarsi delle emozioni. Sa esattamente
quanto valgono le persone che usa — e le usa lo stesso.

\textbf{Nota storica:} La sua effettiva presenza fisica a Palermo nel
1282 è dibattuta. Alcune fonti lo collocano in Sicilia; altre in Aragona.
In gioco, questa ambiguità è produttiva: i personaggi non sono mai
certi di dove sia o cosa stia facendo davvero.

\filetto

\subsection{Jean de Saint-Rémy}
\textit{Il giustiziere angioino di Palermo. Fa il suo lavoro. Questo è il problema.}

\statblock{Jean de Saint-Rémy}{Giustiziere Regio}%
{2}{3}{2}{2}

\textbf{Disponibile a:} Tutti i tavoli (antagonista primario o secondario)

\textbf{Ruolo in gioco:} De Saint-Rémy non è un villain da manuale.
È un funzionario competente che fa il suo lavoro in una situazione
che sfugge al suo controllo. Sa che qualcosa si sta muovendo, ma
non sa dove guardare. È paranoico, reattivo, e disposto a usare
la violenza — ma preferisce la prevenzione alla repressione.

\textbf{Carattere:} Professionale, stressato, convinto di star
difendendo l'ordine legittimo. Non si considera crudele. Si sbaglia.

\textbf{Meccanica speciale — L'orologio di De Saint-Rémy:} Il GM
coordinatore tiene un orologio segreto da 6 segmenti. Ogni volta che
un tavolo compie un'azione che attira l'attenzione delle autorità,
l'orologio avanza. Quando è pieno, De Saint-Rémy ordina uno stato
di emergenza: coprifuoco, chiusura del porto, pattuglie rinforzate.
Questo riduce automaticamente di 1 il punteggio di ogni fazione.

\filetto

\subsection{L'Arcivescovo di Palermo}
\textit{Tra Dio, Roma, e il suo gregge. Tre padroni. Nessuno soddisfatto.}

\statblock{Arcivescovo di Palermo}{Prelato — Guglielmo di Monforte}%
{2}{3}{1}{4}

\textbf{Disponibile a:} Clero (presenza diretta), Baroni (negoziazione),
Congiurati (contatto indiretto)

\textbf{Ruolo in gioco:} L'arcivescovo sa più di quanto mostri.
Ha ricevuto comunicazioni da Roma e da fonti siciliane che lo mettono
in una posizione impossibile. Vuole che il suo popolo stia bene.
Non vuole essere complice di un assassinio di massa. Queste due
cose potrebbero essere incompatibili.

\textbf{Carattere:} Diplomatico, cauto, genuinamente preoccupato.
Non è un eroe — è un uomo che cerca di sopravvivere moralmente
a una situazione in cui sopravvivere moralmente è quasi impossibile.

\secsep

%% ============================================================
%% NPC PER FAZIONE: CONGIURATI
%% ============================================================
\section{NPC dei Congiurati}

\subsection{Emissario di Pietro d'Aragona}
{\cinzel\footnotesize\color{grigionote} Guillem de Montcada}\\
\textit{Parla a nome del re. Non dice tutto quello che sa.}

\statblock{L\'\ Emissario}{Agente Aragonese — Guillem de Montcada}%
{2}{3}{2}{2}

\textbf{Ruolo:} Collegamento tra la rete siciliana e la corte di Barcellona.
Porta istruzioni, oro, e informazioni — non tutte esatte. Pietro d'Aragona
ha i suoi tempi, che non coincidono necessariamente con quelli della
rivolta. L'emissario è il punto dove questa frizione diventa concreta.

\textbf{Segreto:} Sa che Pietro potrebbe non intervenire militarmente
se la rivolta scoppia nei tempi sbagliati. Non ha istruzioni su
come gestire questa informazione.

\subsection{Il Contatto di Messina}
\textit{La Sicilia orientale è un mondo diverso. Ma può fare la differenza.}

\statblock{Contatto di Messina}{Mercante-Agente — Enrico di Taormina}%
{2}{2}{2}{3}

\textbf{Ruolo:} Tiene i fili della rete nella Sicilia orientale.
Se i Congiurati riescono a mantenerlo operativo, la rivolta può
estendersi a Messina — raddoppiando l'impatto. Se la rete viene
compromessa, sparisce.

\subsection{Agente Bizantino}
{\cinzel\footnotesize\color{grigionote} Alexios Philanthropenos}\\
\textit{L'oro di Costantinopoli ha un prezzo. Quale, non è chiaro.}

\statblock{Agente Bizantino}{Diplomatico-Spia — Alexios Philanthropenos}%
{1}{4}{1}{3}

\textbf{Ruolo:} Rappresenta gli interessi di Michele VIII Paleologo.
Finanzia la cospirazione ma vuole garanzie che la spedizione di Carlo
venga bloccata — indipendentemente da ciò che accade alla Sicilia.
I suoi interessi e quelli dei Congiurati si sovrappongono senza
coincidere.

\secsep

%% ============================================================
%% NPC PER FAZIONE: BARONI
%% ============================================================
\section{NPC dei Baroni}

\subsection{Alaimo da Lentini}
\textit{Il barone più potente della Sicilia orientale. E il più prudente.}

\statblock{Alaimo da Lentini}{Barone}%
{3}{3}{2}{3}

\textbf{Disponibile a:} Baroni (figura di riferimento), Congiurati
(alleato potenziale)

\textbf{Ruolo in gioco:} Alaimo è il peso che potrebbe spostare
l'equilibrio — ma non si muove facilmente. Ha visto troppi piani
fallire per fidarsi del primo che arriva. I personaggi dei Baroni
possono tentare di convincerlo, ma devono farlo con argomenti concreti.
Un Alaimo convinto vale automaticamente +1 al punteggio della fazione.

\textbf{Carattere:} Pragmatico, lento, capace di lealtà assoluta
verso chi guadagna la sua fiducia. Non perdona i tradimenti.

\textbf{Nota storica:} Alaimo da Lentini è storicamente attestato
come uno dei leader della rivolta nella Sicilia orientale. Fu poi
implicato in una contro-cospirazione contro Pietro d'Aragona e
giustiziato nel 1287. Il suo arco narrativo reale è una tragedia.

\filetto

\subsection{Il Vassallo Fedele agli Angioini}
\textit{Ha scelto Carlo. Non se ne è ancora pentito.}

\statblock{Vassallo Filo-Angioino}{Barone Minore — Ugo di Santapau}%
{2}{2}{2}{2}

\textbf{Ruolo:} Antagonista interno alla fazione. Non è un nemico
dichiarato, ma le sue scelte informano i francesi — consapevolmente
o no. I personaggi dei Baroni devono decidere cosa farne: neutralizzarlo,
convertirlo, o ignorarlo.

\subsection{La Famiglia Rivale}
\textit{Non vi odiate per via di Carlo. Vi odiate per via del 1260.}

\statblock{Capofamiglia Rivale}{Barone — Ruggero di Incisa}%
{2}{2}{3}{2}

\textbf{Ruolo:} Una faida di vecchia data complica l'unità baronale.
I personaggi devono gestire questa tensione senza che esploda —
o usarla strategicamente. Un'alleanza impossibile, se realizzata,
vale +1 al punteggio.

\secsep

%% ============================================================
%% NPC PER FAZIONE: CLERO
%% ============================================================
\section{NPC del Clero}

\subsection{Fra' Matteo}
{\cinzel\footnotesize\color{grigionote} Fra' Matteo di Sant'Agata}\\
\textit{Frate minore. Sa tutto quello che succede nei quartieri poveri.}

\statblock{Fra' Matteo di Sant'Agata}{Frate Minore}%
{3}{2}{2}{3}

\textbf{Ruolo:} Informatore naturale per il Clero. Circola nei
quartieri dell'Albergheria e della Kalsa, conosce le famiglie,
sa dove sono i soldati francesi la sera. Non è un cospiratore —
è una persona che vuole che il suo popolo stia bene, e usa quello
che sa per aiutarlo.

\textbf{Carattere:} Allegro in superficie, profondamente stanco
sotto. Ha visto troppe cose che avrebbe preferito non vedere.

\filetto

\subsection{Il Legato Pontificio}
\textit{Inviato di Martino IV. Filo-angioino. Ma non stupido.}

\statblock{Legato Pontificio}{Emissario Papale — Gherardo da Parma}%
{1}{4}{1}{4}

\textbf{Ruolo:} Antagonista complesso per il Clero. Non è nemico
della Sicilia — è fedele a Roma, che in questo momento coincide
con Carlo. Se i personaggi del Clero riescono a fargli capire
cosa sta davvero succedendo, potrebbe scegliere di non intervenire.
Non aiuterà — ma potrebbe non ostacolare.

\subsection{La Monaca Visionaria}
\textit{Ha avuto una visione. La interpreta. È pericoloso.}

\statblock{Suor Eufemia di Corleone}{Monaca}%
{2}{2}{1}{4}

\textbf{Ruolo:} Figura ambigua che può stabilizzare o destabilizzare.
Le sue visioni vengono prese sul serio dalla popolazione. Se il
Clero riesce a inquadrarle narrativamente (le visioni parlano di
liberazione), diventano uno strumento potente. Se perdono il
controllo, possono scatenare una reazione incontrollata.

\secsep

%% ============================================================
%% NPC PER FAZIONE: POPOLO
%% ============================================================
\section{NPC del Popolo}

\subsection{Il Soldato Francese}
\textit{Ha vent'anni. È lontano da casa. Ha paura anche lui.}

\statblock{Soldato Francese}{Milite Angioino — Renaud de Coucy}%
{2}{1}{3}{1}

\textbf{Ruolo:} Antagonista umano. Non è un mostro. È un giovane
che esegue ordini in un paese che non capisce. I personaggi del
Popolo possono scegliere di odiarlo, di compatirlo, o di usarlo.
Ognuna di queste scelte ha conseguenze diverse.

\textbf{Nota di design:} Questo NPC è deliberatamente costruito
per creare disagio morale. La violenza contro di lui deve avere
un costo emotivo.

\filetto

\subsection{Il Capoquartiere}
\textit{Tiene l'ordine nell'Albergheria. Da decenni.}

\statblock{Capoquartiere}{Figura Comunitaria — Salamone dell'Albergheria}%
{3}{2}{2}{3}

\textbf{Ruolo:} Alleato potenziale del Popolo. Non è un rivoluzionario —
è un uomo pratico che vuole che il suo quartiere sopravviva. Se i
personaggi lo convincono che la rivolta è inevitabile, organizzerà
il quartiere. Se non lo convincono, si metterà al riparo.

\subsection{La Madre del Morto}
{\cinzel\footnotesize\color{grigionote} Addolorata di Seralcadio}\\
\textit{Suo figlio è stato ucciso dai francesi tre mesi fa. Nessuno
ha risposto. Lei non ha dimenticato.}

\statblock{La Madre del Morto}{Testimone — Addolorata di Seralcadio}%
{1}{2}{1}{4}

\textbf{Ruolo:} Non è un'agente — è una testimonianza vivente.
Se i personaggi riescono a portarla in pubblico a raccontare la
sua storia, il tiro di Aizzare successivo è automaticamente Eccezionale.
Il suo dolore è reale. Usarlo è una scelta morale che i personaggi
devono fare consapevolmente.

\filettodoppio
